\section{Manus：上下文工程是生产 Agent 的核心}

Kimi、ChatGPT Agent、Qwen 分别展示模型、产品和代码 RL，本章看生产 Agent 最容易被低估的一层：上下文工程。Manus 选择基于前沿模型的 in-context 能力构建 Agent，而不是从头训练端到端模型。它强调上下文工程是一门实验科学，生产 Agent 的关键不只是 prompt，而是 KV cache、工具管理、文件系统记忆、todo 复述、保留错误和避免少样本陷阱。

\begin{figure}[H]
\centering
\includegraphics[width=0.96\textwidth]{manus-context-engineering.png}
\caption{Manus Context Engineering：KV cache、工具遮蔽、文件系统、todo、错误保留和多样性。自制概念图，依据 Manus 官方博文和 01:59:04--02:06:06 对谈内容整理。}
\end{figure}

\begin{knowledgebox}{读图：上下文工程是运行时工程}
训练让模型有能力，上下文工程让 Agent 在生产中稳定、便宜、可恢复地使用能力。它处理的是每一轮工具调用、上下文增长和失败恢复。
\end{knowledgebox}

\subsection{KV cache 与工具遮蔽}

上一节说上下文工程是运行时工程，本节拆两个最实用的原则：KV cache 和工具遮蔽。Manus 认为 KV-cache 命中率是生产 Agent 最重要指标之一，因为 Agent 的输入输出比例高度偏向输入，缓存直接影响延迟和成本。它还主张“遮蔽，而非移除”工具：动态增删工具会破坏缓存并混淆模型；更好的做法是保留工具定义，用状态机或 logits mask 控制可选工具。

\begin{knowledgebox}{术语消化：上下文工程关键点}
\begin{tabularx}{\textwidth}{p{0.22\textwidth}p{0.34\textwidth}X}
\toprule
技术 & 作用 & 错误做法 \\
\midrule
KV cache & 降低长前缀重复推理成本 & 每轮改变系统 prompt 或工具定义。 \\
Tool masking & 控制当前可用动作 & 动态删除工具导致上下文不一致。 \\
Filesystem as context & 用文件系统做外部长期记忆 & 把所有内容塞进上下文窗口。 \\
Todo.md & 通过复述目标操控注意力 & 长任务中不刷新目标。 \\
Keep errors & 保留失败证据帮助恢复 & 清理错误导致重复犯错。 \\
\bottomrule
\end{tabularx}
\end{knowledgebox}

\begin{importantbox}{生产 Agent 检查清单}
一个生产 Agent 至少要回答六个问题：上下文是否可缓存？工具空间是否可控？长期记忆在哪里？失败是否保留？目标是否会被复述到近期上下文？重复示例是否会让模型陷入模式？Manus 的经验几乎都围绕这些问题展开。
\end{importantbox}

\begin{knowledgebox}{生产 Agent 系统 checklist}
\begin{tabularx}{\textwidth}{p{0.22\textwidth}p{0.34\textwidth}X}
\toprule
检查项 & 问题 & 不做的后果 \\
\midrule
缓存 & 前缀是否稳定、上下文是否 append-only？ & 成本和延迟暴涨。 \\
工具 & 工具定义是否稳定、动作是否可约束？ & 选错工具、缓存失效、动作幻觉。 \\
记忆 & 长期状态是否外部化到文件或数据库？ & 上下文爆炸或信息丢失。 \\
错误 & 失败轨迹是否保留给模型学习？ & 重复同一错误。 \\
安全 & 高风险动作是否确认和审计？ & Agent 误操作造成真实损害。 \\
\bottomrule
\end{tabularx}
\end{knowledgebox}

\subsection{本章小结}

Manus 提醒我们，Agent 的生产问题常常不是模型不够强，而是上下文、工具、记忆和失败恢复没有设计好。

